\DocumentMetadata{
  pdfversion=2.0,
  pdfstandard=ua-2,
  lang=en-US,
  tagging=on,
  tagging-setup={math/setup=mathml-SE,tabsorder=structure}
}
\documentclass[11pt,letterpaper]{article}
\usepackage[margin=0.68in,includefoot]{geometry}
\usepackage{newcomputermodern}
\usepackage{amsmath,amscd,mathtools}
\usepackage{tikz,adjustbox}
\providecommand{\square}{\mdlgwhtsquare}
\usepackage[hidelinks]{hyperref}
\hypersetup{
  pdftitle={Complex Analysis PhD Exam, May 2000},
  pdfauthor={University of Florida Department of Mathematics},
  pdfsubject={Graduate mathematics examination},
  pdfdisplaydoctitle=true,
  bookmarksopen=true
}
\setcounter{secnumdepth}{0}
\setlength{\parindent}{0pt}
\setlength{\parskip}{0.28em}
\setlength{\emergencystretch}{3em}
\setlength{\leftmargini}{2.15em}
\setlength{\leftmarginii}{2.65em}
\setlength{\leftmarginiii}{3em}
\renewcommand{\labelenumi}{\textbf{\theenumi.}}
\pagestyle{empty}
\raggedbottom
\allowdisplaybreaks
\newenvironment{examproblems}[1]{%
  \begin{enumerate}%
  \setcounter{enumi}{#1}%
  \setlength{\topsep}{0.35em}%
  \setlength{\partopsep}{0pt}%
  \setlength{\itemsep}{0.48em}%
  \setlength{\parsep}{0.18em}%
}{\end{enumerate}}
\newenvironment{examsubparts}{%
  \begin{enumerate}%
  \setlength{\topsep}{0.18em}%
  \setlength{\partopsep}{0pt}%
  \setlength{\itemsep}{0.16em}%
  \setlength{\parsep}{0.1em}%
}{\end{enumerate}}
\newenvironment{examinstructions}{%
  \par\begingroup\itshape
}{%
  \par\endgroup\vspace{0.18em}
}
\makeatletter
\renewcommand\section{\@startsection{section}{1}{0pt}%
  {-1.25ex plus -.25ex minus -.1ex}%
  {0.55ex plus .1ex}%
  {\normalfont\large\bfseries}}
\renewcommand{\@maketitle}{%
  \newpage\null\vskip -1em%
  {\footnotesize\bfseries
    \href{https://www.ufl.edu/}{University of Florida}, \href{https://math.ufl.edu/}{Department of Mathematics}\par}%
  \vskip -0.5em%
  \tagstructbegin{tag=Title}%
    {\LARGE\bfseries\href{https://gma.math.ufl.edu/exams/complex-analysis-phd/}{Complex Analysis PhD Exam}, May 2000\par}%
  \tagstructend%
  \vskip 0.25em%
  \tagmcbegin{artifact}%
    \hrule height 1pt%
  \tagmcend%
  \vskip 1em%
}
\makeatother
\title{Complex Analysis PhD Exam, May 2000}
\author{}
\date{}
\begin{document}
\maketitle
\thispagestyle{empty}
\begin{examinstructions}
Notes. Please write up solutions to the eight problems 1–8, below. Please write LARGE, as the grader’s eyes are older and weaker than your eyes.

\par
A square \(S\subset\mathbb C\) is a set of the form \([a,a+L]\times[b,b+L]\), where~\(L\) is positive. Use “\(s:=\mathrm{Foo}\)” to mean that Foo is the definition of the new symbol~\(s\).
\end{examinstructions}
\begin{examproblems}{0}
\item
Find complex numbers \(a,b,c,d\), with \(ad-bc\ne0\), so that the Möbius transformation 
\[
\mu(z):=\frac{az+b}{cz+d}
\]
 carries the imaginary axis to the circle whose radius is~\(2\) and whose center is \(3=3+0i\).
\item
With~\(B\) the open unit ball \(|z|<1\), consider a non-constant analytic function \(h:B\to\mathbb C\).
\begin{examsubparts}
\item[\textnormal{i:}]
Suppose that \(\Re(h(z))\ge0\) for each \(z\in B\). Prove that the inequality can then be strengthened to “\(>\)”.
\item[\textnormal{ii:}]
With \(\Re(h(z))>0\) on~\(B\), suppose further that \(h(0)=1\). Prove, for each \(z\in B\), that 
\[
\frac{1-|z|}{1+|z|}\le |h(z)|\le\frac{1+|z|}{1-|z|}.
\]
\end{examsubparts}
\item
Suppose that \(P()\) is a monic polynomial with degree \(N\ge1\). Let \(\alpha_1,\ldots,\alpha_N\) be an enumeration (with multiplicity) of the zeros of \(P()\).
\begin{examsubparts}
\item[\textnormal{①:}]
Suppose that \(\forall k:\ \Re(\alpha_k)>0\). Prove that all the zeros of the derivative, \(P'\), also lie in the positive half-plane, as follows: Establish that 
\[
\frac{P'(z)}{P(z)}=\frac{1}{z-\alpha_1}+\frac{1}{z-\alpha_2}+\cdots+\frac{1}{z-\alpha_N},
\]
 then use it to complete the proof.
\item[\textnormal{②:}]
Prove Lucas’s theorem: If all the zeros of a non-constant polynomial~\(P\) lie in a convex polygon \(Q\subset\mathbb C\), then all the zeros of \(P'\) lie in~\(Q\).
\item[\textnormal{③:}]
Show that ② can fail if \(P()\) is allowed to be a rational function: Namely, by letting 
\[
P(z):=\frac{z}{z^2+1},
\]
 find a half-plane~\(H\) which owns a zero of \(P'\) but has no zero of~\(P\).
\end{examsubparts}
\item
Use the Residue Calculus to compute 
\[
I:=\int_0^{+\infty}\frac{1}{[x^4+4]\cdot[x^2+9]^9}\,dx.
\]
 To save arithmetic, you may define some explicit points \(P_1,\ldots,P_L\in\mathbb C\) (what should~\(L\) be?) and explicit functions \(h_1,\ldots,h_L\), and then may express your answer in the form 
\[
I=[h_1(P_1)+\cdots+h_L(P_L)]\cdot\text{Constant}.
\]
 (Do not bother to perform the function-evaluation.)
\item
This problem has two parts.
\begin{examsubparts}
\item[\textnormal{a:}]
State (but do not prove) Morera’s Theorem. (You may use this without proof in (b), if you wish.)
\item[\textnormal{b:}]
Prove this version of the Schwarz Reflection Principle: Suppose~\(f\) is continuous in the closed upper half-plane \(H=\mathbb R\times[0,\infty)\) and is analytic on the interior of~\(H\). Further suppose that~\(f\) is real-valued on the real-axis. By defining \(\Phi=f\) on~\(H\), and 
\[
\Phi(z):=\overline{f(\overline z)},\qquad\text{for all }z\in\mathbb C\setminus H,
\]
 extend~\(f\) to all of \(\mathbb C\). Then this~\(\Phi\) is analytic.
\end{examsubparts}
\item
This problem has three parts.
\begin{examsubparts}
\item[\textnormal{①:}]
Please state Picard’s Theorem.
\item[\textnormal{②:}]
Let~\(f\) be meromorphic in the whole complex plane. Suppose that the range of~\(f\) omits three distinct values (one of them can be~\(\infty\)). Prove that~\(f\) is constant.
\item[\textnormal{③:}]
Suppose that~\(f\) and~\(g\) are entire functions such that, on \(\mathbb C\), 
\[
f^3+g^3=1.
\]
 Prove that~\(f\) and~\(g\) are each constant functions. [Note: Symbol \(f^3\) means \(f\cdot f\cdot f\).]
\end{examsubparts}
\item
Fix a real \(b>0\). Write down an entire function, \(f\), that vanishes precisely on the sequence \((z_n)_{n=1}^{\infty}\), where \(z_n:=n^b\).
\item
Show that all roots of polynomial \(P(z):=z^5+15z+1\) lie in the ball \(|z|<2\), but that only one root satisfies \(|z|<\frac32\).
\end{examproblems}
\end{document}
